% ------------------------------------------------------------------------
%
\begin{frame}{Binary trees}

  \begin{itemize}

    \item A \textbf{binary tree} is a special case of tree, where each
      node has either no children or exactly two.
      \begin{center}
        \includegraphics{bt_shape}
      \end{center}

    \item Nodes are of two kinds: \textbf{internal}
      ({\tiny\raisebox{1pt}{\(\oslash\)}}, \(\circ\) and \(\bullet\))
      or \textbf{external} (\(\scriptstyle \Box\)).

    \item The \textbf{root} is the only internal node without a parent
      ({\tiny\raisebox{1pt}{\(\oslash\)}}).

    \item \textbf{Leaves} (\(\bullet\)) are internal nodes whose
      children are two external nodes (\(\scriptstyle \Box\)).

  \end{itemize}

\end{frame}

% ------------------------------------------------------------------------
%
\begin{frame}{Binary trees}

  Internal nodes are usually associated with some kind of information,
  whilst external nodes are not, like
  \begin{center}
    \includegraphics[bb=71 645 190 721]{bt_ex1}
    \qquad
    \includegraphics[bb=71 645 190 721]{bt_ex2}
  \end{center}
  \centerline{Extended binary tree \qquad\quad Pruned binary tree}

  We need to ponder how many shapes a binary tree can take, just like
  stacks. Each shape as a value constructor and maybe some data.
  \begin{itemize}

    \item The \textbf{empty tree} is an external node, without
      information: \(\fun{ext}().\)

    \item A \textbf{non\hyp{}empty tree} is an internal node carrying
      some information \(x\), and which has a left child~\(t_1\) and a
      right child~\(t_2\): \(\fun{int}(x, t_1, t_2).\)

  \end{itemize}

\end{frame}

% ------------------------------------------------------------------------
%
\begin{frame}{Binary trees}

  The previous tree in the following term:
  \begin{tabbing}
    $\fun{int}(8,$ \= $\fun{int}(1,$ \= $\fun{int}(3,$ \= $\fun{ext}(),$\\
                   \>                \>                \> $\fun{int}(5, \fun{ext}(), \fun{ext}()),$\\
                   \>                \> $\fun{int}(9, \fun{ext}(), \fun{ext}()),$\\
                   \> $\fun{int}(3,$ \= $\fun{ext}(),$\\
                   \>                \> $\fun{int}(2, \fun{ext}(), \fun{ext}())).$
  \end{tabbing}
  We may want to define a function that sums the integers in a binary
  tree like so:
  \begin{equation*}
    \begin{array}{r@{\;}l@{\;}l@{}}
      \fun{add}(\fun{ext}()) & \rightarrow & 0;\\
      \fun{add}(\fun{int}(x, left, right)) & \rightarrow & x + \fun{add} \; left + \fun{add} \; right.
    \end{array}
  \end{equation*}

  The only issue is that we wanted a function that does \emph{not}
  apply to empty trees (it is partial), but we need the value~\(0\)
  for the recursion to work.

\end{frame}

% ------------------------------------------------------------------------
%
\begin{frame}{Optional values}

  We need another function that checks whether the \textbf{whole tree}
  is empty, and, if, returns a special value that means `No sum is
  possible'.

  \bigskip

  The solution is to use a data type called the \textbf{option type}:
  We simply use $\fun{none}()$ to denote that no value is intended,
  and `$\fun{some} \; v$' to mean that the value~$v$ is here.

  \bigskip

  All we need is to call another function, \fun{add'} that handles the
  optional (final) value, like so:
  \begin{equation*}
    \begin{array}{r@{\;}l@{\;}l@{}}
      \fun{add'}(\fun{ext}()) & \rightarrow & \fun{none}();\\
      \fun{add'}(\fun{int} \; node) & \rightarrow & \fun{some} (\fun{add} \; node).
    \end{array}
  \end{equation*}
  Notice how we treated $\fun{int}(x, left, right)$ as a tuple, which
  we named \emph{node}. To make it easier to know that a function is
  undefined (it is used to denote data structures), we will have its
  name start with a capital letter, e.g. `\fun{None}' and `\fun{Int}'.

\end{frame}
